% =====================================================================
% Слайд 3 — Что технически делает система
% =====================================================================
\begin{frame}{Что технически делает система}
  \small
  Задача формализуется как \rlhi{score-following}: по поступающему
  MIDI/аудио-потоку и заранее известной партитуре в реальном времени
  определять текущую позицию исполнителя.

  \begin{columns}[T,onlytextwidth]
    \column{0.42\textwidth}
    \begin{itemize}\setlength{\itemsep}{0.4em}
      \item \textbf{Вход:} партитура (известна) + поток нот солиста
            (realtime).
      \item \textbf{Выход:} индекс ноты $\hat\imath_t$
            $+$ темповый коэффициент $\tau_t$.
      \item \textbf{Бюджет задержки:} $\le 20$ мс
            (порог восприятия рассинхронизации, Cont 2010).
    \end{itemize}

    \column{0.58\textwidth}
    \centering
    % СХЕМА 2 — Pipeline системы (вход → Core ML → выход + RL на боку)
\begin{tikzpicture}[
  every node/.style={font=\footnotesize, align=center},
  box/.style={draw, thick, minimum height=0.9cm, minimum width=1.5cm,
              rounded corners=1.5pt, fill=white, inner sep=2pt},
  small/.style={font=\scriptsize},
  arr/.style={-{Latex[length=2mm]}, thick},
]
  \node[box]                            (in)   {MIDI /\\микрофон};
  \node[box, right=0.4cm of in]         (sf)   {Score\\follower};
  \node[box, right=0.4cm of sf]         (tt)   {Tempo\\Tracker};
  \node[box, right=0.4cm of tt]         (rend) {Renderer};
  \node[box, right=0.4cm of rend]       (out)  {Звуковая\\карта};

  \draw[arr] (in)   -- (sf);
  \draw[arr] (sf)   -- node[above, small] {$\hat\imath_t$} (tt);
  \draw[arr] (tt)   -- node[above, small] {$\tau_t$} (rend);
  \draw[arr] (rend) -- (out);

  % Bracket: Core ML over sf + tt
  \draw[decorate, decoration={brace, amplitude=4pt, raise=2pt}]
        ($(sf.north west) + (0,0.05)$) --
        ($(tt.north east) + (0,0.05)$)
        node[midway, above=8pt, small] {\itshape Core ML};

  % RL agent (dashed, off to the side)
  \node[box, dashed, below=1.1cm of tt, fill=black!4] (rl)
        {\bfseries RL-агент\\\scriptsize\itshape (план)};
  \draw[arr, dashed] (rl) -- (tt);
  \node[small, right=1pt of rl, anchor=west, text width=4cm]
        {\itshape опережающее\\предсказание темпа\\на 200--500 мс};
\end{tikzpicture}

    \figcaption{Core ML (HMM/HSMM/OLTW) и TempoTracker --- классическое
             ядро. RL-агент \emph{дополняет} ядро опережающим
             предсказанием темпа на 200--500 мс.}
  \end{columns}
\end{frame}

